\documentclass[11pt]{article}

\usepackage[a4paper,landscape,margin=0.6in]{geometry}
\usepackage{amsmath,amssymb}
\usepackage{array}
\usepackage{booktabs}
\usepackage{longtable}
\usepackage{hyperref}

\newcommand{\I}{\mathcal{I}}
\newcommand{\DeltaI}{\Delta^{\I}}
\newcommand{\sem}[1]{\left(#1\right)^{\I}}
\newcommand{\Nat}{\mathbb{N}}
\newenvironment{compacttable}
  {\begingroup\setlength{\tabcolsep}{3pt}}
  {\endgroup}

\title{Runir's Description Logic Feature Syntax and Semantics}
\author{Dominik Drexler\\\href{mailto:dominik.drexler@liu.se}{dominik.drexler@liu.se}}
\date{Version 0.0.23}

\begin{document}
\maketitle

\section{Preliminaries}

Let $\I$ be an interpretation with non-empty domain $\DeltaI$.
Each atomic concept $A$ is interpreted as a set
$\sem{A} \subseteq \DeltaI$, each atomic role $R$ as a binary relation
$\sem{R} \subseteq \DeltaI \times \DeltaI$, and each individual name $a$
as an element $\sem{a} \in \DeltaI$.

Concept descriptions denote subsets of $\DeltaI$. Role descriptions denote
binary relations over $\DeltaI$.
For planning states, let $s$ denote the set of atoms true in the current state,
and let $\gamma$ denote the set of goal literals.

\section{Concept Constructors}

\begin{compacttable}
\begin{longtable}{>{\raggedright\arraybackslash}p{0.26\textwidth}
                  >{\centering\arraybackslash}p{0.15\textwidth}
                  >{\raggedright\arraybackslash}p{0.55\textwidth}}
\toprule
Name & Syntax & Semantics \\
\midrule
\endhead

Top &
$\top$ &
$\sem{\top} = \DeltaI$ \\

Bottom &
$\bot$ &
$\sem{\bot} = \emptyset$ \\

Atomic state concept &
$P$ &
$\sem{P} = \{ a \in \DeltaI \mid P(a) \in s \}$ \\

Positive atomic goal concept &
$G^+$ &
$\sem{{G^+}} = \{ a \in \DeltaI \mid G(a) \in \gamma \}$ \\

Negative atomic goal concept &
$G^-$ &
$\sem{{G^-}} = \{ a \in \DeltaI \mid \neg G(a) \in \gamma \}$ \\

Intersection &
$C \sqcap D$ &
$\sem{C \sqcap D} = \sem{C} \cap \sem{D}$ \\

Union &
$C \sqcup D$ &
$\sem{C \sqcup D} = \sem{C} \cup \sem{D}$ \\

Negation &
$\neg C$ &
$\sem{\neg C} = \DeltaI \setminus \sem{C}$ \\

Value restriction &
$\forall R.C$ &
$\sem{\forall R.C} = \{ a \in \DeltaI \mid \forall b.\,(a,b)\in\sem{R} \Rightarrow b\in\sem{C} \}$ \\

Existential quantification &
$\exists R.C$ &
$\sem{\exists R.C} = \{ a \in \DeltaI \mid \exists b.\,(a,b)\in\sem{R} \land b\in\sem{C} \}$ \\

At-least number restriction &
$\geq n\,R$ &
$\sem{\geq n\,R} = \{ a \in \DeltaI \mid |\{ b \in \DeltaI \mid (a,b)\in\sem{R} \}| \geq n \}$ \\

At-most number restriction &
$\leq n\,R$ &
$\sem{\leq n\,R} = \{ a \in \DeltaI \mid |\{ b \in \DeltaI \mid (a,b)\in\sem{R} \}| \leq n \}$ \\

Exact number restriction &
$= n\,R$ &
$\sem{= n\,R} = \{ a \in \DeltaI \mid |\{ b \in \DeltaI \mid (a,b)\in\sem{R} \}| = n \}$ \\

Qualified at-least number restriction &
$\geq n\,R.C$ &
$\sem{\geq n\,R.C} = \{ a \in \DeltaI \mid |\{ b \in \DeltaI \mid (a,b)\in\sem{R} \land b\in\sem{C} \}| \geq n \}$ \\

Qualified at-most number restriction &
$\leq n\,R.C$ &
$\sem{\leq n\,R.C} = \{ a \in \DeltaI \mid |\{ b \in \DeltaI \mid (a,b)\in\sem{R} \land b\in\sem{C} \}| \leq n \}$ \\

Qualified exact number restriction &
$= n\,R.C$ &
$\sem{= n\,R.C} = \{ a \in \DeltaI \mid |\{ b \in \DeltaI \mid (a,b)\in\sem{R} \land b\in\sem{C} \}| = n \}$ \\

Role-value map inclusion &
$R \subseteq S$ &
$\sem{R \subseteq S} = \{ a \in \DeltaI \mid \forall b.\,(a,b)\in\sem{R} \Rightarrow (a,b)\in\sem{S} \}$ \\

Agreement &
$R \doteq S$ &
$\sem{R \doteq S} = \{ a \in \DeltaI \mid \forall b.\,(a,b)\in\sem{R} \Leftrightarrow (a,b)\in\sem{S} \}$ \\

Role fillers &
$R:\{a_1,\ldots,a_k\}$ &
$\sem{R:\{a_1,\ldots,a_k\}} = \{ a \in \DeltaI \mid \{\sem{a_1},\ldots,\sem{a_k}\} \subseteq \{ b \in \DeltaI \mid (a,b)\in\sem{R} \} \}$ \\

One-of &
$\{a_1,\ldots,a_k\}$ &
$\sem{\{a_1,\ldots,a_k\}} = \{\sem{a_1},\ldots,\sem{a_k}\}$ \\

Nominal &
$\{a\}$ &
$\sem{\{a\}} = \{\sem{a}\}$ \\

\bottomrule
\end{longtable}
\end{compacttable}

Here $n \in \Nat$. In role fillers and one-of constructors, the listed
object names are interpreted as individual names in $\I$.

\section{Concrete Syntax for Concept Constructors}

\begin{compacttable}
\begin{longtable}{>{\raggedright\arraybackslash}p{0.30\textwidth}
                  >{\raggedright\arraybackslash}p{0.46\textwidth}
                  >{\raggedright\arraybackslash}p{0.20\textwidth}}
\toprule
Name & Concrete syntax & Abstract syntax \\
\midrule
\endhead
Top & \texttt{c\_top} & $\top$ \\
Bottom & \texttt{c\_bot} & $\bot$ \\
Atomic state concept & \texttt{(c\_atomic\_state "p")} & $P$ \\
Positive atomic goal concept & \texttt{(c\_atomic\_goal "p" true)} & $G^+$ \\
Negative atomic goal concept & \texttt{(c\_atomic\_goal "p" false)} & $G^-$ \\
Intersection & \texttt{(c\_and C1 ... Cn)} & $C_1 \sqcap \cdots \sqcap C_n$ \\
Union & \texttt{(c\_or C1 ... Cn)} & $C_1 \sqcup \cdots \sqcup C_n$ \\
Negation & \texttt{(c\_not C)} & $\neg C$ \\
Value restriction & \texttt{(c\_all R C)} & $\forall R.C$ \\
Existential quantification & \texttt{(c\_some R C)} & $\exists R.C$ \\
At-least number restriction & \texttt{(c\_at\_least n R)} & $\geq n\,R$ \\
At-most number restriction & \texttt{(c\_at\_most n R)} & $\leq n\,R$ \\
Exact number restriction & \texttt{(c\_exactly n R)} & $= n\,R$ \\
Qualified at-least restriction & \texttt{(c\_at\_least n R C)} & $\geq n\,R.C$ \\
Qualified at-most restriction & \texttt{(c\_at\_most n R C)} & $\leq n\,R.C$ \\
Qualified exact restriction & \texttt{(c\_exactly n R C)} & $= n\,R.C$ \\
Same-as / agreement & \texttt{(c\_same\_as R1 R2)} & $R_1 \doteq R_2$ \\
Role-value map & \texttt{(c\_subset R1 R2)} & $R_1 \subseteq R_2$ \\
Role fillers & \texttt{(c\_fillers R a1 ... an)} & $R:\{a_1,\ldots,a_n\}$ \\
One-of & \texttt{(c\_one\_of a1 ... an)} & $\{a_1,\ldots,a_n\}$ \\
Nominal & \texttt{(c\_nominal a)} & $\{a\}$ \\
\bottomrule
\end{longtable}
\end{compacttable}

\paragraph{Polarity of atomic state vs.\ goal constructors.}
Atomic \emph{state} concepts and roles take no polarity argument: the concrete
syntax is \texttt{(c\_atomic\_state "p")} and \texttt{(r\_atomic\_state "p")},
and a negated state concept/role is built compositionally with
\texttt{c\_not}/\texttt{r\_complement}. Atomic \emph{goal} concepts and roles
carry an explicit polarity, \texttt{(c\_atomic\_goal "p" true)} /
\texttt{(c\_atomic\_goal "p" false)} (and likewise for \texttt{r\_atomic\_goal}),
because the goal condition may assert a literal positively or negatively. The
boolean atomic constructors below also carry a polarity argument and require a
nullary (0-arity) predicate.

\section{Restrictions on Role Interpretations}

\paragraph{Functional roles.}
A feature $f$ is interpreted as a functional binary relation:
\[
  \forall a,b,c.\,
  (a,b)\in\sem{f} \land (a,c)\in\sem{f}
  \Rightarrow b=c.
\]
Equivalently, features may be viewed as partial functions.

\paragraph{Transitive roles.}
A transitive role $R$ is interpreted as a transitive binary relation:
\[
  \forall a,b,c.\,
  (a,b)\in\sem{R} \land (b,c)\in\sem{R}
  \Rightarrow (a,c)\in\sem{R}.
\]

\section{Role Constructors}

\begin{compacttable}
\begin{longtable}{>{\raggedright\arraybackslash}p{0.26\textwidth}
                  >{\centering\arraybackslash}p{0.15\textwidth}
                  >{\raggedright\arraybackslash}p{0.55\textwidth}}
\toprule
Name & Syntax & Semantics \\
\midrule
\endhead

Universal role &
$U$ &
$\sem{U} = \DeltaI \times \DeltaI$ \\

Atomic state role &
$P$ &
$\sem{P} = \{ (a,b) \in \DeltaI \times \DeltaI \mid P(a,b) \in s \}$ \\

Positive atomic goal role &
$G^+$ &
$\sem{{G^+}} = \{ (a,b) \in \DeltaI \times \DeltaI \mid G(a,b) \in \gamma \}$ \\

Negative atomic goal role &
$G^-$ &
$\sem{{G^-}} = \{ (a,b) \in \DeltaI \times \DeltaI \mid \neg G(a,b) \in \gamma \}$ \\

Intersection &
$R \sqcap S$ &
$\sem{R \sqcap S} = \sem{R} \cap \sem{S}$ \\

Union &
$R \sqcup S$ &
$\sem{R \sqcup S} = \sem{R} \cup \sem{S}$ \\

Complement &
$\neg R$ &
$\sem{\neg R} = (\DeltaI \times \DeltaI) \setminus \sem{R}$ \\

Inverse &
$R^{-}$ &
$\sem{R^{-}} = \{(b,a)\in\DeltaI\times\DeltaI \mid (a,b)\in\sem{R}\}$ \\

Composition &
$R \circ S$ &
$\sem{R \circ S} = \{(a,c) \mid \exists b.\,(a,b)\in\sem{R} \land (b,c)\in\sem{S}\}$ \\

Transitive closure &
$R^{+}$ &
$\sem{R^{+}} = \bigcup_{n\geq 1}(\sem{R})^n$ \\

Reflexive-transitive closure &
$R^{*}$ &
$\sem{R^{*}} = \bigcup_{n\geq 0}(\sem{R})^n$ \\

Role restriction &
$R|C$ &
$\sem{R|C} = \sem{R} \cap (\DeltaI \times \sem{C})$ \\

Identity &
$\mathrm{id}(C)$ &
$\sem{\mathrm{id}(C)} = \{(d,d) \mid d\in\sem{C}\}$ \\

\bottomrule
\end{longtable}
\end{compacttable}

The iterated composition $(\sem{R})^n$ is defined by
\[
  (\sem{R})^0 = \{(d,d) \mid d\in\DeltaI\},
  \qquad
  (\sem{R})^{n+1} = (\sem{R})^n \circ \sem{R}.
\]

\section{Concrete Syntax for Role Constructors}

\begin{compacttable}
\begin{longtable}{>{\raggedright\arraybackslash}p{0.30\textwidth}
                  >{\raggedright\arraybackslash}p{0.46\textwidth}
                  >{\raggedright\arraybackslash}p{0.20\textwidth}}
\toprule
Name & Concrete syntax & Abstract syntax \\
\midrule
\endhead
Universal role & \texttt{r\_universal} & $U$ \\
Atomic state role & \texttt{(r\_atomic\_state "p")} & $P$ \\
Positive atomic goal role & \texttt{(r\_atomic\_goal "p" true)} & $G^+$ \\
Negative atomic goal role & \texttt{(r\_atomic\_goal "p" false)} & $G^-$ \\
Intersection & \texttt{(r\_and R1 ... Rn)} & $R_1 \sqcap \cdots \sqcap R_n$ \\
Union & \texttt{(r\_or R1 ... Rn)} & $R_1 \sqcup \cdots \sqcup R_n$ \\
Complement & \texttt{(r\_complement R)} & $\neg R$ \\
Inverse & \texttt{(r\_inverse R)} & $R^{-}$ \\
Composition & \texttt{(r\_composition R1 ... Rn)} & $R_1 \circ \cdots \circ R_n$ \\
Transitive closure & \texttt{(r\_transitive\_closure R)} & $R^{+}$ \\
Reflexive-transitive closure & \texttt{(r\_reflexive\_transitive\_closure R)} & $R^{*}$ \\
Role restriction & \texttt{(r\_restriction R C)} & $R|C$ \\
Identity & \texttt{(r\_identity C)} & $\mathrm{id}(C)$ \\
\bottomrule
\end{longtable}
\end{compacttable}

\section{Boolean Constructors}

\begin{compacttable}
\begin{longtable}{>{\raggedright\arraybackslash}p{0.30\textwidth}
                  >{\centering\arraybackslash}p{0.20\textwidth}
                  >{\raggedright\arraybackslash}p{0.46\textwidth}}
\toprule
Name & Syntax & Semantics \\
\midrule
\endhead

Positive atomic state predicate &
$\mathrm{holds}(p)$ &
$\mathrm{true}$ iff there exists an atom $p() \in s$ \\

Negative atomic state predicate &
$\neg\mathrm{holds}(p)$ &
$\mathrm{true}$ iff no current atom with predicate $p$ is true in the state \\

Positive atomic goal predicate &
$\mathrm{goal}(p)$ &
$\mathrm{true}$ iff $p() \in \gamma$ \\

Negative atomic goal predicate &
$\neg\mathrm{goal}(p)$ &
$\mathrm{true}$ iff $\neg p() \in \gamma$ \\

Non-empty concept &
$C \neq \emptyset$ &
$\mathrm{true}$ iff $\sem{C} \neq \emptyset$ \\

Non-empty role &
$R \neq \emptyset$ &
$\mathrm{true}$ iff $\sem{R} \neq \emptyset$ \\

\bottomrule
\end{longtable}
\end{compacttable}

\section{Concrete Syntax for Boolean Constructors}

\begin{compacttable}
\begin{longtable}{>{\raggedright\arraybackslash}p{0.30\textwidth}
                  >{\raggedright\arraybackslash}p{0.46\textwidth}
                  >{\raggedright\arraybackslash}p{0.20\textwidth}}
\toprule
Name & Concrete syntax & Abstract syntax \\
\midrule
\endhead
Positive atomic state predicate & \texttt{(b\_atomic\_state "p" true)} & $\mathrm{holds}(p)$ \\
Negative atomic state predicate & \texttt{(b\_atomic\_state "p" false)} & $\neg\mathrm{holds}(p)$ \\
Positive atomic goal predicate & \texttt{(b\_atomic\_goal "p" true)} & $\mathrm{goal}(p)$ \\
Negative atomic goal predicate & \texttt{(b\_atomic\_goal "p" false)} & $\neg\mathrm{goal}(p)$ \\
Non-empty concept & \texttt{(b\_nonempty C)} & $C \neq \emptyset$ \\
Non-empty role & \texttt{(b\_nonempty R)} & $R \neq \emptyset$ \\
\bottomrule
\end{longtable}
\end{compacttable}

\section{Numerical Constructors}

\begin{compacttable}
\begin{longtable}{>{\raggedright\arraybackslash}p{0.30\textwidth}
                  >{\centering\arraybackslash}p{0.20\textwidth}
                  >{\raggedright\arraybackslash}p{0.46\textwidth}}
\toprule
Name & Syntax & Semantics \\
\midrule
\endhead

Concept count &
$|C|$ &
$|\sem{C}|$ \\

Role count &
$|R|$ &
$|\sem{R}|$ \\

Role distance &
$d_R(C,D)$ &
$\min\{n \in \Nat \mid \sem{C} \times \sem{D} \cap (\sem{R})^n \neq \emptyset\}$, or $\infty$ if empty \\

\bottomrule
\end{longtable}
\end{compacttable}

\section{Concrete Syntax for Numerical Constructors}

\begin{compacttable}
\begin{longtable}{>{\raggedright\arraybackslash}p{0.30\textwidth}
                  >{\raggedright\arraybackslash}p{0.46\textwidth}
                  >{\raggedright\arraybackslash}p{0.20\textwidth}}
\toprule
Name & Concrete syntax & Abstract syntax \\
\midrule
\endhead
Concept count & \texttt{(n\_count C)} & $|C|$ \\
Role count & \texttt{(n\_count R)} & $|R|$ \\
Role distance & \texttt{(n\_distance C R D)} & $d_R(C,D)$ \\
\bottomrule
\end{longtable}
\end{compacttable}


\section{Unsolvability Family Constructors}

The unsolvability family (\texttt{uns}) extends the base family with comparison operators and
constants. A comparison takes two operands of the same category---either two Booleans or two
Numericals---and yields a Boolean. Boolean operands are interpreted numerically as $0$ (false) or
$1$ (true), so the same relational operators apply uniformly to both operand categories. We write
$\sem{e}$ for the numeric value of an operand $e$ (a count/distance for Numericals,
or $0/1$ for Booleans).

\begin{compacttable}
\begin{longtable}{>{\raggedright\arraybackslash}p{0.30\textwidth}
                  >{\centering\arraybackslash}p{0.20\textwidth}
                  >{\raggedright\arraybackslash}p{0.46\textwidth}}
\toprule
Name & Syntax & Semantics \\
\midrule
\endhead

Equal &
$e_1 = e_2$ &
$\mathrm{true}$ iff $\sem{e_1} = \sem{e_2}$ \\

Not equal &
$e_1 \neq e_2$ &
$\mathrm{true}$ iff $\sem{e_1} \neq \sem{e_2}$ \\

Less than &
$e_1 < e_2$ &
$\mathrm{true}$ iff $\sem{e_1} < \sem{e_2}$ \\

Less or equal &
$e_1 \leq e_2$ &
$\mathrm{true}$ iff $\sem{e_1} \leq \sem{e_2}$ \\

Greater than &
$e_1 > e_2$ &
$\mathrm{true}$ iff $\sem{e_1} > \sem{e_2}$ \\

Greater or equal &
$e_1 \geq e_2$ &
$\mathrm{true}$ iff $\sem{e_1} \geq \sem{e_2}$ \\

Boolean constant &
$b \in \{\mathrm{true}, \mathrm{false}\}$ &
the constant truth value $b$ \\

Numerical constant &
$k \in \Nat$ &
the constant nonnegative integer $k$ \\

\bottomrule
\end{longtable}
\end{compacttable}

\section{Concrete Syntax for Unsolvability Family Constructors}

\begin{compacttable}
\begin{longtable}{>{\raggedright\arraybackslash}p{0.30\textwidth}
                  >{\raggedright\arraybackslash}p{0.46\textwidth}
                  >{\raggedright\arraybackslash}p{0.20\textwidth}}
\toprule
Name & Concrete syntax & Abstract syntax \\
\midrule
\endhead
Boolean equal & \texttt{(b\_eq B1 B2)} & $B_1 = B_2$ \\
Boolean not equal & \texttt{(b\_neq B1 B2)} & $B_1 \neq B_2$ \\
Boolean less than & \texttt{(b\_lt B1 B2)} & $B_1 < B_2$ \\
Boolean less or equal & \texttt{(b\_le B1 B2)} & $B_1 \leq B_2$ \\
Boolean greater than & \texttt{(b\_gt B1 B2)} & $B_1 > B_2$ \\
Boolean greater or equal & \texttt{(b\_ge B1 B2)} & $B_1 \geq B_2$ \\
Numerical equal & \texttt{(n\_eq N1 N2)} & $N_1 = N_2$ \\
Numerical not equal & \texttt{(n\_neq N1 N2)} & $N_1 \neq N_2$ \\
Numerical less than & \texttt{(n\_lt N1 N2)} & $N_1 < N_2$ \\
Numerical less or equal & \texttt{(n\_le N1 N2)} & $N_1 \leq N_2$ \\
Numerical greater than & \texttt{(n\_gt N1 N2)} & $N_1 > N_2$ \\
Numerical greater or equal & \texttt{(n\_ge N1 N2)} & $N_1 \geq N_2$ \\
Boolean constant & \texttt{(b\_const true)} / \texttt{(b\_const false)} & $\mathrm{true}$ / $\mathrm{false}$ \\
Numerical constant & \texttt{(n\_const k)} & $k$ \\
\bottomrule
\end{longtable}
\end{compacttable}


\section{Module Programs}

A module program state is an extended state
\[
  x = (\bar{s}, m, \mu, \rho_C, \rho_R)
\]
where $\bar{s}$ is an annotated planning state, $m$ is the current module,
$\mu$ is a typed memory state, and $\rho_C,\rho_R$ are the concept and role
register valuations.  The implementation keeps memory as
$\mu \in M_{\mathit{int}} \uplus M_{\mathit{ext}}$ so internal control states
and externally visible module states are not conflated.

A proof edge records the transition that changed the planning state and the
rule that justified the transition.  For sketch proofs the edge label is
$(e,r)$ with a state-graph edge $e$ and a sketch rule $r$.  For module-program
proofs the corresponding label is $(e,r_v)$ where $r_v$ is a rule variant
(load, do, sketch, or call).  Edges that only expose control failure or an open
state may omit $e$ or $r_v$ until the parser/executor rejects the corresponding
ill-formed program earlier.

The intended small-step relation has the form
\[
  (\bar{s},m,\mu,\rho_C,\rho_R) \xrightarrow{r_v,e}
  (\bar{s}',m',\mu',\rho'_C,\rho'_R).
\]
Load rules update registers and internal memory, do and sketch rules may change
$\bar{s}$ through a planning transition, and call rules push the caller frame
before entering the callee module.  A proof succeeds when the initial extended
state reaches a goal-annotated state and the constructed proof graph has no
open states and no directed cycle.

\section*{Source}

This document summarizes Appendix~1, Section~A1.2 of
\emph{The Description Logic Handbook}, edited by Franz Baader,
Deborah L. McGuinness, Daniele Nardi, and Peter F. Patel-Schneider.

\end{document}
